% Compile from the repository root with XeLaTeX, without shell escape.
% Clarify conjugate transposes in the existing proof; the original PDF
% already uses overlined transposes, whose bars are lost in text extraction.
% Match the original/reference: 11 pt Fandol Song, Computer Modern math,
% 164 mm text width and 1.12 line spacing. Only page 2's proof opening changes.
\documentclass[11pt]{article}
\usepackage{amsmath,amssymb}
\usepackage[no-math]{fontspec}
\usepackage{xeCJK}
\setmainfont{Latin Modern Roman}
\setCJKmainfont{FandolSong-Regular.otf}[Path=assets/fonts/fandol/,BoldFont=FandolSong-Bold.otf,ItalicFont=FandolSong-Regular.otf,Script=CustomDefault]
\xeCJKsetup{PunctStyle=quanjiao,CheckSingle=true}
\linespread{1.12}
\usepackage[active,tightpage]{preview}
\setlength{\textwidth}{164mm}
\setlength{\PreviewBorder}{4pt}
\PreviewEnvironment{preview}
\begin{document}
\begin{preview}
\begin{minipage}{164mm}
\setlength{\parindent}{0pt}
\setlength{\parskip}{3pt}
本题中 $Z^*=\overline Z^{\mathsf T}$ 表示共轭转置。将(1)两式分别左乘 $x^*,y^*$，相加得
\[
x^*Ax+x^*By+y^*B^{\mathsf T}x+y^*Cy-\mathrm i\,x^*Dx=0.
\]
实对称阵亦为 Hermite 阵，故 $A^*=A,\ C^*=C$，从而 $x^*Ax,y^*Cy\in\mathbb R$。又因 $B$ 为实矩阵，
\[
y^*B^{\mathsf T}x=\overline{x^*By},\qquad
x^*By+y^*B^{\mathsf T}x=2\operatorname{Re}(x^*By)\in\mathbb R.
\]
故前四项之和为实数。因 $D$ 为实对角阵，$x^*Dx\in\mathbb R$，取虚部得
\[
0=-x^*Dx=-\sum_{j=1}^n d_j|x_j|^2.
\]
\end{minipage}
\end{preview}
\end{document}
